\scanpage{5}
\mspage{3}
On \unsure{posera|pose.}
\[
\underline{\deleted{\Phi}\mathrm{Ch}}_\infty(X)=\underline{\mathrm{Hom}}(\underset{\inserted{\text{procatégorie}}}{\mathbb{\Phi}},X)\overset{\mathrm{déf}}{=}
\deleted{\unsure{\varinjlim_n\underline{\Phi}_n(X)}}\inserted{\varinjlim\underline{\Phi}_n(X)}
\]
\inserted{ou simplement $\underline{\mathrm{Ch}}(X)$}
% slanted
\margin{L'ens des objets de cette catégorie se notera $\mathrm{Ch}_\infty(X)$ ou $\mathrm{Ch}(X)$, de \unsure{$=$} on définit les $\mathrm{Ch}_n(X)$}

On a une inclusion pleinement fidèle
\[
\underline{\deleted{\Phi}\mathrm{Ch}}_\infty(X)\lhook\joinrel\longrightarrow\underline{\mathrm{Hom}}(\Phi,X),
\]
où la deuxième catégorie est celle des \deleted{suites} diagrammes infinis (dans les deux directions)
\diagram{\begin{tikzcd}[ampersand replacement=\&,column sep=.85em]
\cdots \& A_{-2}\arrow[l]\arrow[r,"u_{-2}"] \& A_{-1} \& \deleted{\unsure{X}}A_0\arrow[l,"u_{-1}"']\arrow[r,"u_0"] \& \deleted{\unsure{X}}A_1 \& A_2\arrow[l,"u_1"']\arrow[r,"u_2"] \& \cdots
\end{tikzcd}}
et où $\underline{\mathrm{Ch}}_\infty(X)$ s'identifie à la sous-catégorie pleine formée des objets pour lesquels les flèches $u_i$ sont des identités, pour $|i|$ assez grand.

On appellera $\underline{\mathrm{Ch}}_\infty(X)$ la \uline{catégorie des chemins} \unsure{dans} $X$, ses objets sont les chemins de $X$. Ceux qui proviennent de $\Phi_n(X)$ seront appelés \uline{chemins \unsure{standart|standard} de longueur $2n$}. (Le terme ``standart'' s'applique ici à nos conventions de notations. Ainsi nous ne considérons que des chemins de longueur \unsure{pair|paire}, et pour une longueur déterminée, nous choisissons \inserted{l'une} parmi les deux \deleted{possibilités de les décrire} \inserted{espèces possibles}, p.ex. pour les chemins de longueur $2$ nous avons pris seulement
\diagram{\begin{tikzcd}[ampersand replacement=\&]
A_{-1} \& A_0\arrow[l]\arrow[r] \& A_1
\end{tikzcd}}
et non pas
\diagram{\begin{tikzcd}[ampersand replacement=\&]
A_{-1}\arrow[r] \& A_0 \& A_1\arrow[l] \& .
\end{tikzcd}}

% archive sheet number: 4
